\raggedbottom % evita que LaTeX estire el espacio vertical cuando una figura pasa a la siguiente página
aggedbottom % evita que LaTeX estire el espacio vertical cuando una figura pasa a la siguiente página
\pagenumbering{arabic} %necessary to start numbering pages in arabic numerals from here
\setcounter{page}{1}

\section{Introducción}

Una proteína es una cadena de aminoácidos, y la forma tridimensional que adopta en el agua determina casi todo lo que puede hacer: qué moléculas reconoce, qué reacciones cataliza y con qué otras proteínas se une. Conocer la secuencia de aminoácidos es relativamente sencillo; conocer la forma a la que se pliega es mucho más difícil. Predecir esa forma a partir de la secuencia es uno de los problemas centrales de la bioinformática, y durante décadas fue también uno de los más estudiados en computación, porque el número de formas posibles crece de manera exponencial con la longitud de la cadena.

Para estudiar el problema sin toda la complejidad química de una proteína real, Dill propuso en 1985 el modelo hidrofóbico-polar, o modelo HP~\cite{dill1985theory, lau1989lattice}. En este modelo cada aminoácido es solo de dos tipos, hidrofóbico (H) o polar (P), la cadena se dibuja sobre una cuadrícula y la energía de una forma depende únicamente de cuántos pares de aminoácidos H quedan juntos. A pesar de ser tan simple, encontrar la forma de menor energía en el modelo HP es un problema NP-completo~\cite{crescenzi1998complexity, berger1998protein}, y se ha convertido en un problema de prueba clásico para algoritmos de búsqueda.

En esta práctica aplicamos la Optimización por Enjambre de Partículas (PSO) al modelo HP en dos dimensiones. A diferencia de las prácticas anteriores, donde las soluciones eran vectores de números reales, aquí una solución es una forma de la cadena en la cuadrícula, que es un objeto discreto y que además no puede cruzarse consigo mismo. Por eso la parte más importante del trabajo fue decidir cómo representar un plegamiento con los números continuos que maneja PSO. Elegimos una representación con direcciones relativas, en la que cada partícula guarda tres valores de preferencia por cada paso de la cadena (girar a la izquierda, seguir de frente o girar a la derecha), y un decodificador que repara los choques de la cadena consigo misma.

También hicimos variar los coeficientes $w$, $c_1$ y $c_2$ durante la ejecución, pero en lugar de una variación lineal usamos una curva coseno a la que se suma un recalentamiento cuando el mejor valor del enjambre deja de mejorar. Probamos el algoritmo en cinco cadenas de referencia de la literatura, de 20 a 64 aminoácidos, con 30 ejecuciones independientes por cadena. En cada iteración registramos el mejor valor encontrado (\textit{gbest}) y tres medidas de diversidad del enjambre (de posición, de velocidad y cognitiva), y reportamos su media y desviación estándar. Por último, comparamos la estrategia de coeficientes elegida contra las curvas lineal, exponencial y coseno sin recalentamiento con la prueba de Mann-Whitney.
